% TEMPLATE-MILC-v3.tex - plantilla maestra UNICA de las guias MILC v3.
%
% La usan las tres familias de guias, cada una con su driver:
%   · Grados 6-11  -> scripts/build-guias-g11.py
%   · Bebras       -> scripts/build-guias-bebras.py
%   · Semillero    -> scripts/build-guias-semillero.py
%
% Antes habia tres copias casi identicas (~400 lineas iguales, 6 distintas):
% cada mejora habia que aplicarla tres veces, y en la practica no se hacia
% (el motor de recursos vivia solo en dos de ellas). Lo que varia entre
% familias esta parametrizado en 6 placeholders que cada driver llena
% (se nombran aqui SIN su sintaxis real: el builder reemplaza por texto plano,
% tambien dentro de los comentarios, y un valor multilinea rompe el preambulo):
%
%   HEADER_IZQ        encabezado izquierdo de pagina
%   HEADER_DER        encabezado derecho de pagina
%   PORTADA_ETIQUETA  rotulo sobre el numero grande (GRADO/MOMENTO/MODULO)
%   PORTADA_SERIE     linea de serie de la portada
%   PORTADA_ID        identificador de la guia en la portada
%   PORTADA_CIRCULO   el circulito de grado (°), o vacio si no aplica
%
% El resto de claves las documenta content/guias/_SCHEMA.md. El script valida
% que no queden placeholders sin reemplazar antes de escribir el archivo final.

\documentclass[11pt,letterpaper]{article}
\usepackage[margin=2.0cm]{geometry}
\usepackage[table]{xcolor}
\usepackage{fontspec}
\usepackage[spanish]{babel}
\usepackage{tikz}
% Librerías TikZ para los diagramas de las guías (bloque `recursos:`):
%  · babel  → babel en español vuelve activos caracteres como «>», y TikZ los usa
%             en sus opciones (>=stealth, ->). Sin esta librería, cualquier
%             diagrama con flechas revienta con «\language@active@arg> has an
%             extra }». Debe ir ANTES de que se dibuje nada.
%  · positioning        → sintaxis «below left=2pt and 3pt».
%  · decorations.pathreplacing → llaves ({brace}) para señalar rangos.
\usetikzlibrary{babel,positioning,decorations.pathreplacing}
\usepackage{graphicx}
% Circuitos: el trabajo de electrónica se hace hoy con micro:bit y sus sensores
% integrados (MakeCode, sin cableado), así que no se necesitan esquemáticos.
% Si algún día entran montajes con protoboard, basta descomentar esta línea:
% los diagramas .tex/.tikz declarados en `recursos:` ya se incrustan solos.
% \usepackage{circuitikz}
\usepackage[most]{tcolorbox}
\usepackage{tabularx}
\usepackage{array}
\usepackage{fancyhdr}
\usepackage{titlesec}
\usepackage[document]{ragged2e}  % bandera a la derecha en todo el cuerpo
\usepackage{enumitem}
\usepackage{microtype}

% Helvetica Neue no trae la flecha U+2192, asi que cada «→» escrito literalmente
% en el YAML se compilaba a NADA, en silencio (462 flechas en 61 guias: rutas del
% tipo «paleta → tipografia → lockup» salian rotas). Mapeamos el caracter a su
% version matematica, que si tiene glifo. Asi el autor puede seguir escribiendo
% «→» comodamente en el YAML.
\usepackage{newunicodechar}
\newunicodechar{→}{\ensuremath{\rightarrow}}
% Mismo problema con los simbolos de marca y vinetas: el «✓» de «una palomita ✓»
% salia como cuadrito vacio en el PDF de todas las guias que lo usaban.
\usepackage{pifont}
\newunicodechar{✓}{\ding{51}}
\newunicodechar{✗}{\ding{55}}
\newunicodechar{✔}{\ding{52}}
\newunicodechar{•}{\textbullet}
\newunicodechar{≥}{\ensuremath{\geq}}
\newunicodechar{≤}{\ensuremath{\leq}}

\setmainfont{Helvetica Neue}
\setsansfont{Helvetica Neue}
\setlength{\parindent}{0pt}
\setlength{\parskip}{6pt}
% Sin particion de palabras: en 6.º-7.º una palabra cortada por guion al final
% de linea («Comuni-cacion») frena la lectura mucho mas que una linea corta.
\hyphenpenalty=10000
\exhyphenpenalty=10000
\pretolerance=2000
\tolerance=3000
\setlength{\emergencystretch}{3em}
\linespread{1.10}
\renewcommand{\arraystretch}{1.25}
\setlist{leftmargin=5mm,itemsep=1mm,topsep=1mm}
\newcolumntype{Y}{>{\RaggedRight\arraybackslash}X}

\definecolor{milcMagenta}{HTML}{E6007E}
\definecolor{milcVino}{HTML}{5A0038}
\definecolor{milcTurquesa}{HTML}{0093A5}
\definecolor{milcVerde}{HTML}{58A923}
\definecolor{milcMostaza}{HTML}{E5B400}
\definecolor{milcOcre}{HTML}{B86600}
\definecolor{milcNegro}{HTML}{241816}
\definecolor{milcGris}{HTML}{F7F7F4}
\definecolor{milcLinea}{HTML}{E8E2DD}
\definecolor{milcGradePrimary}{HTML}{FF2D87}
\definecolor{milcGradeDark}{HTML}{9F1B56}
\definecolor{milcGradeSoft}{HTML}{FFE0EE}
\definecolor{milcGradeAccent}{HTML}{CFE7FF}
\definecolor{milcGradeText}{HTML}{FFFFFF}
\color{milcNegro}

\pagestyle{fancy}
\fancyhf{}
\lhead{\small Tecnología e Informática · Grado 10}
\rhead{\small Guía 21 · MILC v3}
\cfoot{\small\thepage}
\renewcommand{\headrulewidth}{0pt}

\newtcolorbox{titlebox}[1]{
  enhanced,colback=#1,colframe=#1,arc=7mm,boxrule=0pt,
  left=7mm,right=7mm,top=3mm,bottom=3mm,
  before skip=8pt,after skip=8pt,
  fontupper=\sffamily\bfseries\Large\color{white}
}

\newtcolorbox{softbox}[3]{
  enhanced,colback=#2,colframe=#1,borderline west={4pt}{0pt}{#1},
  arc=2mm,boxrule=.4pt,left=4mm,right=4mm,top=3mm,bottom=3mm,
  before skip=6pt,after skip=8pt,title={#3},coltitle=white,
  colbacktitle=#1,fonttitle=\sffamily\bfseries,fontupper=\RaggedRight,
  attach boxed title to top left={xshift=3mm,yshift=-2mm},
  boxed title style={arc=2mm, boxrule=0pt}
}

\newtcolorbox{drawbox}[2]{
  enhanced,colback=white,colframe=#1,arc=4mm,boxrule=1pt,
  left=5mm,right=5mm,top=4mm,bottom=4mm,before skip=8pt,after skip=8pt,
  title={#2},coltitle=white,colbacktitle=#1,fonttitle=\sffamily\bfseries,
  fontupper=\RaggedRight,
  attach boxed title to top left={xshift=4mm,yshift=-2mm},
  boxed title style={arc=3mm, boxrule=0pt}
}

\newtcolorbox{infoband}[1]{
  enhanced,colback=#1!8,colframe=#1!50,arc=2mm,boxrule=.5pt,
  left=4mm,right=4mm,top=2mm,bottom=2mm,before skip=4pt,after skip=4pt,
  fontupper=\small\RaggedRight
}

% Puente narrativo entre secciones (contrato editorial Conéctate).
% Da continuidad: "Acabas de X. Ahora vas a Y porque Z."
% Visualmente: banda izquierda en color del grado, texto en itálica pequeña.
\newtcolorbox{puentebox}{
  enhanced,colback=milcGradeSoft,colframe=milcGradeDark,
  borderline west={3pt}{0pt}{milcGradeDark},
  arc=0mm,boxrule=0pt,left=5mm,right=4mm,top=2.5mm,bottom=2.5mm,
  before skip=4pt,after skip=8pt,
  fontupper=\itshape\small\color{milcNegro}\RaggedRight
}
% La flecha va en modo matematico a proposito: Helvetica Neue NO tiene el glifo
% U+2192, asi que \textrightarrow se compilaba a NADA («Missing character: There
% is no → in font Helvetica Neue») y el puente salia sin flecha. En modo
% matematico la flecha viene de la fuente matematica y si se dibuja.
\newcommand{\puente}[1]{%
  \begin{puentebox}%
    \textcolor{milcGradeDark}{$\rightarrow$}\hspace{2mm}#1%
  \end{puentebox}%
}

\newcommand{\linead}[1]{\noindent\color{milcLinea}\rule{#1}{0.5pt}\color{milcNegro}}
\newcommand{\respuesta}[1]{\par\vspace{2mm}\foreach \n in {1,...,#1}{\noindent\color{milcLinea}\rule{\linewidth}{0.5pt}\par\vspace{5mm}}\color{milcNegro}}
\newcommand{\checkbox}{\textcolor{milcGradeDark}{\fbox{\phantom{x}}}\hspace{5pt}}

\newcommand{\phasecard}[4]{
\begin{tcolorbox}[enhanced,colback=#3,colframe=#2,arc=3mm,boxrule=.5pt,height=3.05cm,valign=center,halign=center,left=1.5mm,right=1.5mm,top=2mm,bottom=2mm]
{\sffamily\bfseries\fontsize{22}{24}\selectfont\textcolor{#2}{#1}}\par
{\sffamily\bfseries\small #4}
\end{tcolorbox}
}

% ── Piezas del rediseno 2026 (legibilidad 6.º-11.º) ───────────────────
% Banda de estacion: reemplaza el titlebox macizo. Titulo a la izquierda,
% dato practico (tiempo/modalidad) a la derecha, en una sola linea.
\newcommand{\milcestacion}[2]{%
\begin{tcolorbox}[enhanced,colback=#1,colframe=#1,arc=2mm,boxrule=0pt,
  left=5mm,right=5mm,top=2.6mm,bottom=2.6mm,before skip=11pt,after skip=7pt]
{\sffamily\bfseries\fontsize{15}{18}\selectfont\color{white}#2}
\end{tcolorbox}}

% Tarjeta de paso numerada: sustituye las tablas «Pilar / Aplicacion», que
% eran formato de consulta docente, no de estudiante. El numero va en un
% circulo sobre el borde para que la secuencia se lea de un vistazo.
\newcommand{\milcpaso}[3]{%
\begin{tcolorbox}[enhanced,colback=white,colframe=milcLinea,arc=1.5mm,boxrule=.9pt,
  left=14mm,right=4mm,top=3mm,bottom=3mm,before skip=4pt,after skip=4pt,
  fontupper=\RaggedRight,
  overlay={\node[anchor=north west,circle,fill=#2,text=white,inner sep=0pt,
    minimum size=7.4mm,font=\sffamily\bfseries\small]
    at ([xshift=3.6mm,yshift=-3.2mm]frame.north west) {#1};}]
#3
\end{tcolorbox}}

% Casilla marcable de tamano util con lapiz (la anterior era \fbox{\phantom{x}},
% de alto variable segun el texto que la seguia).
\newcommand{\milcmarca}{\framebox[3.8mm]{\rule{0pt}{3.4mm}}}

% Fila marcable de lista suelta (checks de la estacion 1).
\newcommand{\milccheckrow}[1]{%
\par\vspace{1.8mm}\noindent
\makebox[6.4mm][l]{\raisebox{-.55mm}{\textcolor{milcNegro}{\milcmarca}}}%
\parbox[t]{\dimexpr\linewidth-6.4mm\relax}{\RaggedRight #1}\par}

% Celda marcable dentro de la rubrica (tabla de autoevaluacion). Misma
% sangria francesa, pero pensada para vivir dentro de una columna X.
\newcommand{\milcopcion}[1]{%
\makebox[5.2mm][l]{\raisebox{-.55mm}{\textcolor{milcNegro}{\milcmarca}}}%
\parbox[t]{\dimexpr\linewidth-5.2mm\relax}{\RaggedRight #1}}

% ── Recursos visuales de la guía ──────────────────────────────────────
% Declarados en el bloque `recursos:` del YAML y emitidos por el builder.
% \guiaFigura   → imagen rasterizada o PDF (foto, carta celeste, montaje).
% \guiaDiagrama → fuente TikZ/circuitikz (.tex) incrustada como vector:
%                 es la vía para circuitos y esquemáticos editables.
% #1 = ruta relativa al .tex · #2 = pie (puede ir vacío)
\newcommand{\guiaFigura}[2]{%
\begin{center}
\includegraphics[width=0.88\linewidth,height=0.38\textheight,keepaspectratio]{#1}\\[1.5mm]
{\footnotesize\itshape\textcolor{milcNegro!75}{#2}}
\end{center}
}

\newcommand{\guiaDiagrama}[2]{%
\begin{center}
\input{#1}\\[1.5mm]
{\footnotesize\itshape\textcolor{milcNegro!75}{#2}}
\end{center}
}

\begin{document}
\thispagestyle{empty}

% ── Portada ───────────────────────────────────────────────────────────
% Antes: un rectangulo de color a pagina completa. Se veia bien en pantalla,
% pero en la fotocopiadora del colegio se comia el toner y salia gris sucio.
% Ahora la banda ocupa el tercio superior y el resto va en blanco.
\begin{tikzpicture}[remember picture,overlay]
  \path[fill=milcGradePrimary]
    (current page.north west) rectangle ([yshift=-10.6cm]current page.north east);

  \node[anchor=north west,text=milcGradeText] at ([xshift=2.0cm,yshift=-1.55cm]current page.north west)
    {\sffamily\bfseries\fontsize{12}{14}\selectfont GRADO};
  \node[anchor=north west,text=milcGradeText,opacity=.85] at ([xshift=2.0cm,yshift=-2.35cm]current page.north west)
    {\sffamily\bfseries\fontsize{8.5}{10}\selectfont SERIE GUÍAS · TECNOLOGÍA E INFORMÁTICA};
  \node[anchor=north east,text=milcGradeText,opacity=.85,align=right] at ([xshift=-2.0cm,yshift=-1.55cm]current page.north east)
    {\sffamily\bfseries\fontsize{9}{12}\selectfont I.E. Sor María Juliana\\Cartago, Valle del Cauca};

  \node[anchor=west,text=milcGradeText] (gradonum) at ([xshift=1.85cm,yshift=-7.1cm]current page.north west)
    {\sffamily\bfseries\fontsize{112}{112}\selectfont 10};
  % El circulito de grado (°) solo aplica a las guias de grado («11°»); en
  % Bebras y Semillero el numero grande es un momento/modulo, no un grado.
  % Se posiciona relativo al nodo `gradonum`, no en coordenadas absolutas.
  \draw[milcGradeText,line width=1.6pt]
    ([xshift=2.0mm,yshift=-4.5mm]gradonum.north east) circle (.15cm);

  \node[anchor=west,text=milcGradeText,text width=11.4cm,align=left] at ([xshift=1.5cm,yshift=0.5cm]gradonum.east)
    {
     {\sffamily\bfseries\fontsize{9}{11}\selectfont\MakeUppercase{Período 3 · Ofimática con IA: contabilidad, datos y emprendimiento}}\\[3mm]
     {\sffamily\bfseries\fontsize{21}{25}\selectfont\setlength{\baselineskip}{27pt}Ofimática con IA ---\\[4pt]el copiloto\\[4pt]del oficio cotidiano}\\[3.5mm]
     {\sffamily\bfseries\fontsize{15}{18}\selectfont Guía 21-10-TIC}
    };
\end{tikzpicture}

\vspace*{9.4cm}
\vfill

{\sffamily\bfseries\fontsize{8}{10}\selectfont\textcolor{milcGradeDark}{LO QUE TE LLEVAS HOY}}

\begin{tcolorbox}[enhanced,colback=white,colframe=milcNegro,arc=2mm,boxrule=1.6pt,
  left=5mm,right=5mm,top=4mm,bottom=4mm,before skip=3pt,after skip=12pt,
  fontupper=\RaggedRight\fontsize{11}{15.5}\selectfont]
Lista personalizada de 10 tareas reales de ofimática que haces o harás en tu vida próxima (con la familia, en el colegio, en un primer trabajo, en emprendimiento personal) + para cada tarea, propuesta concreta de cómo la IA puede asistir (no reemplazar) + identificación de qué herramientas gratuitas usarías para cada caso (Google Sheets, Excel Online gratis, LibreOffice, Notion, ChatGPT, Gemini, Bing Copilot integrado en Microsoft 365) + reflexión personal de 5 líneas sobre cuáles 3 herramientas vas a priorizar este periodo
\end{tcolorbox}

\begin{tcolorbox}[enhanced,colback=milcGris,colframe=milcGris,arc=2mm,boxrule=0pt,
  left=5mm,right=5mm,top=3.5mm,bottom=3.5mm,after skip=12pt]
{\sffamily\bfseries\fontsize{8}{10}\selectfont\textcolor{milcNegro!55}{ESTA GUÍA ES DE}}\\[3mm]
\begin{tabularx}{\linewidth}{X p{3.2cm} p{3.2cm}}
\linead{\linewidth} & \linead{3.0cm} & \linead{3.0cm}\\[-1mm]
{\fontsize{8.5}{10}\selectfont\textcolor{milcNegro!55}{Nombre completo}} &
{\fontsize{8.5}{10}\selectfont\textcolor{milcNegro!55}{Grupo}} &
{\fontsize{8.5}{10}\selectfont\textcolor{milcNegro!55}{Fecha}}\\
\end{tabularx}
\end{tcolorbox}

\vfill

\noindent\textcolor{milcLinea}{\rule{\linewidth}{1pt}}\\[2mm]
{\sffamily\bfseries\fontsize{9.5}{12}\selectfont Autor: Dr. Álvaro Cárdenas Orozco}\\[1mm]
{\sffamily\fontsize{8.5}{11}\selectfont\textcolor{milcNegro!55}{Metodología MILC · apertura ancestral · pensamiento computacional · triángulo Dussel–estoicismo–Floridi}}

\newpage

\section*{Apertura: Ofimática con IA --- el copiloto del oficio digital cotidiano}

\begin{softbox}{milcGradeDark}{milcGradeSoft}{Saber ancestral}
En las tiendas y microempresas del centro de Cartago, en las panaderías del barrio Obrero, en las modistas del Quindío y en las pequeñas oficinas de contabilidad de los pueblos cafeteros, hubo durante décadas una práctica que cualquier comerciante experimentado conocía: \textbf{la ofimática del oficio antes del computador}. Aunque no se llamaba así, los tenderos, sastres y contadores manejaban \textbf{herramientas de oficina} con disciplina: (1) \textbf{El cuaderno cuadriculado}: la hoja de cálculo ancestral. Filas para días, columnas para conceptos (\emph{venta-mañana, venta-tarde, gastos, saldo}). (2) \textbf{La calculadora de mano}: el procesador de cálculo. Sumar, restar, sacar porcentaje. (3) \textbf{La radio}: el correo en tiempo real. Conocer precios de mercado, noticias del sector. (4) \textbf{El cuaderno de direcciones}: la base de datos de clientes. Cada uno con teléfono, deuda, fecha de últimas compras. Esas 4 piezas (cuaderno + calculadora + radio + agenda) eran la \textbf{ofimática del oficio del barrio}. Quien las usaba con disciplina llevaba el negocio bien; quien improvisaba, se perdía. La sabiduría era ancestral: \textbf{ningún negocio sostenible se llevaba solo con la memoria}. La IA hoy agrega un \emph{copiloto}: el cuaderno digital ya no hace solo cálculos, ahora sugiere, analiza, propone. Pero la decisión sigue siendo del comerciante.
\end{softbox}

\begin{softbox}{milcTurquesa}{milcGris}{Saber contemporáneo}
La \textbf{ofimática asistida por IA} es la combinación de herramientas tradicionales de oficina (procesador de texto, hoja de cálculo, presentaciones, formularios) con asistencia de IA generativa para acelerar tareas, no para reemplazar el criterio. En 2026, las herramientas gratuitas disponibles incluyen: (1) \textbf{Google Sheets}: hoja de cálculo gratuita con cuenta Gmail. Tiene función \texttt{=GEMINI()} integrada en 2025+ para preguntar a IA dentro de celdas. (2) \textbf{Excel Online}: versión web gratuita de Microsoft Excel, incluye Copilot básico. (3) \textbf{LibreOffice}: paquete completo gratis descargable (Writer + Calc + Impress + Base). (4) \textbf{Notion}: base de datos y notas con IA integrada en plan gratuito limitado. (5) \textbf{ChatGPT, Gemini, Bing Copilot, Claude.ai}: asistentes generales gratuitos para preguntas, análisis, redacción. La regla profesional: \textbf{cada herramienta tiene su lugar; aprender a combinarlas vale más que dominar una sola}. Google Sheets para cálculos colaborativos, Notion para bases de datos personales, LibreOffice si prefieres trabajar offline. La IA acelera todas: \emph{``hazme una fórmula que calcule X''}, \emph{``analiza esta tabla''}, \emph{``escribe un correo basado en este registro''}.
\end{softbox}

\begin{softbox}{milcMostaza}{milcGris}{Pregunta-puente}
\textit{¿Qué sabía el tendero al combinar cuaderno + calculadora + radio + agenda, que el novato olvida cuando intenta llevar microempresa solo con la memoria? ¿Y por qué tener IA como copiloto requiere primero conocer las herramientas que va a copilotar?}
\end{softbox}

\begin{softbox}{milcVerde}{milcGris}{Saber y saber hacer}
Al terminar podrás: (1) \textbf{identificar} las 10 tareas reales de ofimática que dominarán tu vida laboral próxima; (2) \textbf{analizar} cuál herramienta gratuita se ajusta mejor a cada tarea; (3) \textbf{evaluar} cómo la IA puede asistir cada tarea sin reemplazar tu criterio; (4) \textbf{explicar} por qué algunas tareas se benefician más de IA que otras y dónde la phronesis humana sigue siendo crucial.
\end{softbox}



\small
\begin{tabularx}{\textwidth}{>{\bfseries}p{3.0cm}Y}
\rowcolor{milcGradePrimary!12}Autor & Dr. Álvaro Cárdenas Orozco\\
DBA & Aplicación de ofimática asistida por IA a contabilidad básica y emprendimiento (MEN, T\&I)\\
Referentes & Lean Startup · Phronesis financiera · Ética del dato empresarial\\
Producto final & Lista personalizada de 10 tareas reales de ofimática que haces o harás en tu vida próxima (con la familia, en el colegio, en un primer trabajo, en emprendimiento personal) + para cada tarea, propuesta concreta de cómo la IA puede asistir (no reemplazar) + identificación de qué herramientas gratuitas usarías para cada caso (Google Sheets, Excel Online gratis, LibreOffice, Notion, ChatGPT, Gemini, Bing Copilot integrado en Microsoft 365) + reflexión personal de 5 líneas sobre cuáles 3 herramientas vas a priorizar este periodo\\
\end{tabularx}
\normalsize

\puente{Antes de empezar, mira el viaje completo. Cerraste el periodo 2 (textos profesionales con IA). Hoy abre el periodo 3: \textbf{ofimática con IA, contabilidad, emprendimiento}. Las próximas 9 sesiones (contabilidad básica, Sheets + IA, flujo de caja, visualización, encuestas de mercado, Canvas, comunicación empresarial, proyecto integrador, sustentación) construyen sobre el panorama de hoy.}

\milcestacion{milcGradeDark}{Ruta MILC: así vamos a aprender}
\begin{tabularx}{\textwidth}{>{\centering\arraybackslash}X>{\centering\arraybackslash}X>{\centering\arraybackslash}X>{\centering\arraybackslash}X}
\phasecard{1}{milcVerde}{milcGris}{Escucha\\[-1mm]\small Observo, escucho y pregunto.} &
\phasecard{2}{milcTurquesa}{milcGris}{Entender\\[-1mm]\small Sistematizo y ordeno.} &
\phasecard{3}{milcMagenta}{milcGris}{Hacer\\[-1mm]\small Construyo y aplico.} &
\phasecard{4}{milcVino}{milcGris}{Cierre\\[-1mm]\small Evalúo y reflexiono.}
\end{tabularx}


\puente{Antes de teorizar, vas a hacer una \emph{lluvia personal}: anota 10 tareas reales de ofimática que crees harás en los próximos 5 años (en tu casa, en el colegio, en un primer trabajo, en un emprendimiento). Sin filtrar por importancia: solo enumera lo que sospechas.}

\milcestacion{milcVerde}{Estación 1 de 3 · Escucha}

\begin{softbox}{milcVerde}{milcGris}{Escena para escuchar}
\textbf{Actividad 1 · IDENTIFICA --- 10 tareas reales de ofimática} (15 min · individual). Anota 10 tareas concretas de ofimática que crees harás en los próximos 5 años. Reglas: (1) tareas reales, no genéricas. (2) cubre 3 ámbitos: tu casa/familia (3-4 tareas), tu colegio o universidad (3-4 tareas), tu primer trabajo o emprendimiento (3-4 tareas). Ejemplos: \emph{llevar la cuenta de gastos familiares mensuales}, \emph{hacer hoja de cálculo del torneo deportivo del colegio}, \emph{registrar clientes del microemprendimiento familiar}, \emph{calcular precios de servicios para vender}, \emph{hacer factura simple para un cliente}, \emph{crear inventario de productos}, \emph{hacer presentación para clase o trabajo}, \emph{escribir hoja de vida}, \emph{enviar correos profesionales con datos adjuntos}, \emph{analizar resultados de encuesta familiar}.
\end{softbox}

{\sffamily\bfseries\fontsize{8}{10}\selectfont\textcolor{milcNegro!55}{MARCA CUANDO LO HAYAS HECHO}}
\milccheckrow{Listé 10 tareas reales en 3 ámbitos.}
\milccheckrow{Cada tarea es específica, no genérica.}
\milccheckrow{Reconozco que dominaré varias en mi vida próxima.}

\begin{infoband}{milcVerde}


\textbf{Lo que va al cuaderno (Actividad 1 · IDENTIFICA).} Título: «Actividad 1 --- 10 tareas reales». \textbf{Formato:} lista numerada de 10 con ámbito (casa, colegio, trabajo). \textbf{Sabes que terminaste cuando} ves que tu vida próxima requiere ofimática.
\end{infoband}

\puente{Ya tienes 10 candidatas. Ahora vas a entender la \textbf{anatomía de la ofimática con IA}: las 5 herramientas gratuitas principales, cuándo asistir con IA, cuándo solo herramienta tradicional, los 5 errores típicos.}

\milcestacion{milcTurquesa}{Estación 2 de 3 · Entender}

\begin{softbox}{milcTurquesa}{milcGris}{Lo que vas a entender}
Tus 10 tareas son el mapa. Ahora necesitas la \textbf{anatomía de las herramientas gratuitas} y cómo asignar IA a cada caso.
\end{softbox}

\milcpaso{1}{milcTurquesa}{Las \textbf{5 herramientas gratuitas principales} se descomponen así: (1) \textbf{Google Sheets}: hoja de cálculo colaborativa con cuenta Gmail. Mejor para contabilidad personal, control de gastos, análisis simple. La función \texttt{=GEMINI()} permite preguntar a IA dentro de la celda (\emph{``=GEMINI(``analiza la columna B y propone categorías'')''}). (2) \textbf{Excel Online}: gratis con cuenta Microsoft. Incluye Copilot básico que genera fórmulas y analiza datos. (3) \textbf{LibreOffice}: paquete completo descargable, gratis, offline. Writer (texto), Calc (hojas), Impress (presentaciones), Base (bases de datos simples). Sin IA integrada nativa, pero compatible con ChatGPT externo. (4) \textbf{Notion}: notas y bases de datos personales con IA limitada en plan gratuito. Ideal para gestión personal de proyectos. (5) \textbf{Asistentes IA generales}: ChatGPT, Gemini, Bing Copilot, Claude.ai. Útiles para preguntas puntuales sobre fórmulas, análisis, redacción.}
\milcpaso{2}{milcTurquesa}{El patrón profesional \textbf{IA como copiloto, no como piloto}: la IA es buena para sugerir fórmulas, proponer análisis preliminares, generar primeros borradores. La IA es mala para: tomar decisiones finales sobre presupuesto, decidir qué categorías son significativas para tu negocio, evaluar consecuencias humanas de decisiones. La phronesis del oficio consiste en \textbf{usar IA donde acelera trabajo mecánico, conservar criterio humano donde hay consecuencias}.}
\milcpaso{3}{milcTurquesa}{Los \textbf{5 errores típicos} de la ofimática con IA: (a) \textbf{Aceptar fórmulas de IA sin verificar}: \texttt{=GEMINI()} puede sugerir fórmulas mal estructuradas. Siempre probar con casos conocidos. (b) \textbf{Pegar datos personales sensibles en IA pública}: nombres, números de cuenta, datos protegidos no van a ChatGPT. (c) \textbf{Depender de Excel/Sheets con suscripción cara}: hay alternativas gratis igual de buenas. (d) \textbf{No aprender las herramientas base}: pretender que la IA sustituye conocer fórmulas básicas. (e) \textbf{Mezclar herramientas sin propósito}: usar Notion + Sheets + Excel + LibreOffice cuando 2 bastan.}
\milcpaso{4}{milcTurquesa}{Algoritmo: (1) revisar las 10 tareas. (2) para cada una, identificar qué herramienta gratuita conviene. (3) para cada una, evaluar cómo IA puede asistir. (4) priorizar 3 herramientas que dominarás este periodo. (5) escribir reflexión sobre tu mapa personal.}

\begin{drawbox}{milcTurquesa}{Anatomía de la ofimática con IA}
\textbf{Sheets:} colaborativo + GEMINI integrada. \\[1mm] \textbf{Excel Online:} gratuito + Copilot básico. \\[1mm] \textbf{LibreOffice:} offline + paquete completo. \\[1mm] \textbf{Notion:} notas + bases simples. \\[1mm] \textbf{Asistentes IA:} ChatGPT, Gemini, Bing, Claude.
\end{drawbox}

\begin{softbox}{milcMostaza}{milcGris}{Errores comunes y su corrección}
(1) Aceptar fórmulas IA sin verificar. (2) Pegar datos sensibles en IA pública. (3) Pagar suscripciones innecesarias. (4) No aprender herramientas base. (5) Mezclar herramientas sin propósito. (6) Asumir que IA reemplaza criterio. (7) Ignorar atajos de teclado que aceleran todo.
\end{softbox}

\begin{infoband}{milcTurquesa}


\textbf{Lo que va al cuaderno (Actividad 2 · ANALIZA).} Título: «Actividad 2 --- Anatomía ofimática IA». \textbf{Formato:} (a) 5 herramientas con uso típico; (b) IA como copiloto, no piloto; (c) 7 errores con marca. \textbf{Sabes que terminaste cuando} puedes asignar herramienta a cualquier tarea.
\end{infoband}

\puente{Tienes el método. Toca aplicarlo: para cada una de las 10 tareas, propón herramienta gratuita y asistencia de IA. Reflexionas sobre las 3 herramientas que priorizarás este periodo.}

\milcestacion{milcMagenta}{Estación 3 de 3 · Hacer}

\begin{softbox}{milcMagenta}{milcGris}{Cómo vas a construir tu producto}
\textbf{Actividad 3 · EVALÚA + EXPLICA --- 10 tareas + propuesta IA + priorización + reflexión} (35 min · individual). Hora de aplicar. Para cada una de las 10 tareas, propón herramienta + asistencia IA. Prioriza 3 herramientas para este periodo. Reflexionas.
\end{softbox}

\milcpaso{1}{milcMagenta}{Tu mapa se construye en 5 pasos: (1) tomar las 10 tareas de la Actividad 1. (2) para cada una, asignar herramienta gratuita en columna. (3) para cada una, escribir 1-2 líneas sobre cómo IA puede asistir. (4) identificar las 3 herramientas que más necesitas dominar este periodo (las que aparecen en más tareas). (5) reflexión personal de 5 líneas sobre prioridades.}
\milcpaso{2}{milcMagenta}{El patrón \textbf{``herramienta antes que asistente''}: aunque la IA es atractiva, primero debes dominar la herramienta base. No tiene sentido pedirle a Gemini que escriba fórmula compleja en Sheets si no conoces qué hace SUMAR.SI o BUSCARV. La phronesis del oficio consiste en \textbf{conocer la herramienta para evaluar lo que la IA propone}.}
\milcpaso{3}{milcMagenta}{La reflexión final responde 3 preguntas: (a) ¿cuáles 3 herramientas vas a priorizar este periodo? (b) ¿qué tareas de las 10 ya sabes hacer bien y cuáles vas a aprender? (c) ¿cómo combinarás IA con criterio humano según el tipo de tarea?}
\milcpaso{4}{milcMagenta}{Algoritmo: (1) tabla 10 tareas × 3 columnas (Tarea, Herramienta gratuita, Asistencia IA). (2) priorizar 3 herramientas. (3) reflexionar 5 líneas. (4) compartir mapa con compañero (opcional).}

\begin{drawbox}{milcMagenta}{Checklist antes de entregar}
\checkbox 10 tareas reales en 3 ámbitos. \\[1mm] \checkbox Herramienta gratuita asignada a cada tarea. \\[1mm] \checkbox Propuesta de IA asistente a cada tarea. \\[1mm] \checkbox 3 herramientas priorizadas para el periodo. \\[1mm] \checkbox Reflexión de 5 líneas. \\[1mm] \checkbox Mapa coherente con tu vida próxima.
\end{drawbox}

\begin{softbox}{milcMagenta}{milcGris}{Plantilla de trabajo}
\textbf{Tabla 10 tareas.} \textit{Tarea + Herramienta + IA asistente.} \\[1mm] \textbf{3 herramientas priorizadas.} \textit{Las que aparecen más.} \\[1mm] \textbf{Reflexión.} \textit{Qué priorizo y por qué.}
\end{softbox}

\begin{infoband}{milcMagenta}


\textbf{Lo que va al cuaderno (Actividad 3 · EVALÚA + EXPLICA).} Título: «Actividad 3 --- Mi mapa ofimática IA». \textbf{Formato:} tabla 10 tareas + 3 priorizadas + reflexión. \textbf{Sabes que terminaste cuando} sabes con cuáles herramientas vas a trabajar este periodo.
\end{infoband}

\puente{Lo que sigue resume \emph{qué} entregas hoy: lista de 10 tareas + propuesta IA + 3 herramientas priorizadas + reflexión. El criterio principal: que tengas mapa claro de qué herramientas dominar primero según tu propio futuro.}

\milcestacion{milcMostaza}{Producto de la sesión}
\textbf{Lista 10 tareas + herramienta + IA asistente + 3 priorizadas}.
\begin{softbox}{milcMostaza}{milcGris}{Criterios de calidad}
(1) 10 tareas reales en 3 ámbitos. (2) Herramienta gratuita asignada por tarea. (3) Propuesta de asistencia IA por tarea. (4) 3 herramientas priorizadas. (5) Reflexión coherente con vida próxima. (6) Distinción IA copiloto vs piloto.
\end{softbox}

\puente{Antes de cerrar, mira la ofimática desde las cinco dimensiones humanas. Aprender herramientas gratuitas es libertad económica; depender solo de pagas es entregar capital a quien las controla.}

\milcestacion{milcVino}{Evaluación liberadora · cinco dimensiones}

\begin{softbox}{milcVino}{milcGris}{Sentido de la evaluación}
Evaluar no es solo revisar una nota. Es reconocer qué aprendí, qué cambió en mí, qué emociones manejé, cómo me relaciono con la ciudad y la comunidad, y qué le dejaré a quien viene detrás.
\end{softbox}

\textbf{Mi reflexión liberadora en cinco dimensiones.} Completa cada frase iniciada:

\small
\begin{tabularx}{\textwidth}{>{\bfseries\RaggedRight\arraybackslash}p{3.4cm}Y>{\RaggedRight\arraybackslash}p{4.6cm}}
\rowcolor{milcVino}\color{white}Dimensión & \color{white}\bfseries Mi respuesta (frame starter) & \color{white}\bfseries Algo que descubrí\\
1. Personal & Lo que más cambió en mí fue \linead{6.5cm} & \linead{4.2cm}\\
2. Emocional & La emoción más fuerte fue \linead{2.5cm} y la regulé \linead{3cm} & \linead{4.2cm}\\
3. Ciudadana & En democracia esto importa porque \linead{6cm} & \linead{4.2cm}\\
4. Local (Cartago) & En mi barrio sería útil para \linead{6.5cm} & \linead{4.2cm}\\
5. Intergeneracional & A mi abuela/abuelo le diría \linead{6.5cm} & \linead{4.2cm}\\
\end{tabularx}
\normalsize

\begin{infoband}{milcVino}
\textbf{Para la próxima sustentación.} Vuelve a esta tabla en seis meses. Verás que las respuestas cambian — no porque lo aprendido cambie, sino porque tú cambias. La evaluación liberadora es una conversación contigo a través del tiempo.
\end{infoband}


\puente{Y para terminar, tres voces sobre el oficio digital. Dussel sobre las herramientas que liberan al microempresario, Séneca sobre la sobriedad de aprender lo necesario, Floridi sobre la IA como copiloto del oficio profesional.}

\milcestacion{milcGradeDark}{Triángulo de pensamiento · cierre filosófico}

\begin{softbox}{milcVino}{milcGris}{Dussel · lente del nosotros}
\textit{Las herramientas gratuitas liberan al microempresario; las suscripciones caras extraen valor de quien menos puede pagarlas.}

Muchos jóvenes asumen que Microsoft Office o Adobe son las únicas opciones serias. La filosofía de la liberación pide lo opuesto: \textbf{las herramientas gratuitas existen y son tan buenas como las pagas para la mayoría de tareas reales}. El microempresario del barrio que arranca su negocio no necesita pagar Microsoft 365: con Google Workspace gratis o LibreOffice puede llevar contabilidad, hacer presentaciones, gestionar clientes. La phronesis económica del oficio digital consiste en \textbf{conocer las alternativas gratuitas antes de comprometerse con suscripciones}.

\textbf{¿Estoy aprendiendo herramientas gratuitas como acto de libertad económica, o asumo que solo lo pago es bueno?}
\end{softbox}
\linead{\linewidth}\\[1mm]
\linead{\linewidth}

\begin{softbox}{milcMostaza}{milcGris}{Estoicismo · Séneca · cuidado interior}
\textit{Aprende lo necesario con sobriedad; el exceso de herramientas dispersa al que las usa.}

Es tentador querer dominar todas las herramientas disponibles. La sabiduría estoica recomienda lo opuesto: \emph{prioriza 2-3 que cubran el 80 por ciento de tus tareas}. Aprender mal 10 herramientas vale menos que aprender bien 3. La phronesis del oficio consiste en \textbf{elegir tu stack y dominarlo}.

\textbf{¿Estoy priorizando 3 herramientas para dominar, o intento aprender todo al mismo tiempo?}
\end{softbox}
\linead{\linewidth}\\[1mm]
\linead{\linewidth}

\begin{softbox}{milcTurquesa}{milcGris}{Floridi · lente de la infoesfera}
\textit{La IA como copiloto profesional es la nueva ética del oficio en la era de la automatización.}

En la era de la IA generativa, los profesionales que usan IA con criterio aceleran su trabajo sin perder calidad. Los que la usan sin criterio producen errores. Los que la rechazan se quedan atrás. La ética profesional del siglo XXI consiste en \textbf{usar IA como copiloto: ella sugiere, tú decides}. Tu mapa de hoy te entrena en esa práctica.

\textbf{¿Estoy usando la IA como copiloto con criterio, o intentando que sea piloto sin mi intervención?}
\end{softbox}
\linead{\linewidth}\\[1mm]
\linead{\linewidth}

\begin{infoband}{milcNegro}
\textbf{Nota para el docente.} Las tres formulaciones del triángulo son la síntesis del autor de la guía, no citas textuales de los pensadores. Si vas a citarlos en clase, remite a la obra original.
\end{infoband}

\milcestacion{milcVino}{Mi autoevaluación · marca una casilla por fila}

{\small
\setlength{\extrarowheight}{2.2pt}
\begin{tabularx}{\textwidth}{>{\RaggedRight\arraybackslash\bfseries}p{3.0cm}YYY}
\rowcolor{milcVino}\color{white}\bfseries Criterio & \color{white}\bfseries Lo logré & \color{white}\bfseries En proceso & \color{white}\bfseries Necesito apoyo\\[.6mm]
Comprensión del tema & \milcopcion{Explico las ideas con mis palabras.} & \milcopcion{Reconozco las ideas, falta precisión.} & \milcopcion{Necesito repasar lo básico.}\\[.8mm]
\rowcolor{milcGris}Pensamiento computacional & \milcopcion{Usé los pilares al armar mi producto.} & \milcopcion{Usé algunos pilares.} & \milcopcion{Necesito releer la estación 2.}\\[.8mm]
Aplicación y producto & \milcopcion{Mi producto cumple los criterios.} & \milcopcion{Está armado, con ajustes pendientes.} & \milcopcion{Aún está en construcción.}\\[.8mm]
\rowcolor{milcGris}Reflexión 5 dimensiones & \milcopcion{Respondí cada una con honestidad.} & \milcopcion{Respondí algunas con detalle.} & \milcopcion{Mis respuestas son muy generales.}\\[.8mm]
Triángulo de pensamiento & \milcopcion{Conecté las 3 voces con mi trabajo.} & \milcopcion{Conecté al menos una.} & \milcopcion{No las relaciono con mi proceso.}\\[.8mm]
\end{tabularx}}

\vspace{3mm}
\textbf{Hoy entendí que:}\\[1mm]
\linead{\linewidth}\\[1mm]
\linead{\linewidth}

\textbf{Mi compromiso para la próxima sesión es:}\\[1mm]
\linead{\linewidth}\\[1mm]
\linead{\linewidth}

\begin{softbox}{milcMostaza}{milcGris}{Cita de despedida · Marco Aurelio}
\textit{``El obstáculo en el camino se vuelve el camino. Lo que estorba al camino se vuelve camino.''}\\[1mm]
Cada error de hoy, bien observado, se convierte en pieza del camino que pisarás mañana.
\end{softbox}

\small
\begin{tabularx}{\textwidth}{>{\bfseries}p{3.6cm}Y}
\rowcolor{milcGradeDark!12}Nombre completo: & \linead{10cm}\\
Fecha: & \linead{4cm} \quad Curso/grupo: \linead{4cm}\\
Mi firma: & \linead{9cm}\\
\end{tabularx}
\normalsize

\begin{infoband}{milcGradeDark}
\textbf{Para la próxima sesión.} Llega con tu mapa de 10 tareas y las 3 herramientas priorizadas. En la sesión 2 vas a aprender \textbf{contabilidad básica}: ingresos, egresos y balance, llevando cuenta personal o familiar real durante 2 semanas. La ofimática de hoy es el panorama; la contabilidad de mañana es la primera implementación práctica.
\end{infoband}

\end{document}
